\chapter{Implementation}

\section{Development environment and reproducibility}

The pipeline was developed in Python in an isolated virtual environment. Its
dependencies cover HTTP collection, HTML parsing, PDF extraction, JSON and CSV
processing, spreadsheet production, and plotting. LLM inference is provided
locally by Ollama. The project is organised as a collection of independent
scripts rather than as a single monolithic application, so that collection,
cleaning, extraction, evaluation, and analysis can be executed and inspected
separately.

Reproducibility requires recording the Python version, package versions, Ollama
version, model identifier, execution date, prompt version, and inference
parameters. The core workflow does not rely on an external commercial LLM API.
However, reproducibility is not perfect across machines: a model update, changed
source page, or modified prompt can alter outputs. Raw input text, generated
JSON, archived raw LLM answers, and collection indexes are therefore retained as
provenance artefacts.

The project directory separates raw data, cleaned data, intermediate
representations, structured outputs, evaluation material, and analysis outputs.
This organisation prevents a processing step from silently overwriting its
input. It also makes it possible to rerun one stage after a correction without
having to repeat the whole workflow. For example, changes to the taxonomy
normalisation can be applied to existing JSON records without downloading the
corpus again or running new LLM inferences.

\section{Project architecture and data flow}

The implementation follows a file-based architecture. Each source document is
represented by one text file and, after extraction, by one JSON record with the
same logical identity. Dedicated directories store the successive
representations: \texttt{Database} for collected text, \texttt{Cleaned\_Database}
for normalised text, several intermediate directories for the experimental
rule-based pipeline, and \texttt{LLM\_Tagged\_Database} for LLM-generated
records. This one-document-per-file principle simplifies error localisation:
when a value appears incorrect, the corresponding JSON output can be compared
directly with the cleaned text and, where available, the raw model response.

Collection scripts are source-specific because the corpus includes HTML pages,
PDF files, internal action sheets, and research-project records with different
structures. They retrieve the source content, extract or convert it to text,
and maintain an index containing information such as the title, local filename,
and source URL when available. The cleaning scripts then reduce formatting noise
caused by HTML extraction and PDF conversion, including duplicated headers,
fragmented lines, inconsistent whitespace, and filename variations.

The initial rule-based pipeline was implemented as a sequence of modules for
segmentation, section detection, entity extraction, candidate extraction,
heuristic scoring, JSON consolidation, and review-table generation. Each module
writes an intermediate output rather than passing data only in memory to the
next stage. This design was useful for debugging and for understanding the
effect of individual rules, even though the final extraction workflow relies
primarily on constrained LLM extraction.

\section{LLM extraction and output validation}

The LLM module constructs a French prompt for each document and submits it to
the local Ollama chat API. The prompt specifies the required JSON keys, the
expected data types, controlled vocabularies for selected fields, and the
instruction not to invent unsupported information. The model is asked to return
an empty string or an empty list when the document does not provide a reliable
value. The configuration uses the local \texttt{qwen2.5:1.5b} model, a
temperature of 0, a context window of 8,192 tokens, a maximum generation length
of 700 tokens, and an input limit of 2,500 characters for the RFVAA extraction
script.

Structured generation does not guarantee that every response is immediately
usable. The implementation therefore includes defensive processing before a
record is saved. It removes text outside the JSON object when present, repairs
superficial formatting issues such as typographic quotation marks or trailing
commas, and attempts to parse the result. Records that remain invalid can be
flagged for reprocessing or manual inspection rather than being silently merged
into the final dataset. Where available, the unmodified model response is saved
alongside the structured record, which permits later inspection of a parsing
problem or an unexpected extraction.

The resulting JSON records preserve both extracted metadata and contextual
descriptions. They include title, territory, territorial scale, target
audience, initiative type, initiative holder, dates, themes, and source
identification, together with short descriptive fields. The schema is designed
for documentary exploration rather than to claim that every field is known for
every source. Empty values are therefore meaningful: they represent missing or
insufficiently supported information, not a processing failure by default.

\section{Provenance and taxonomy normalisation}

Two deterministic post-processing components complement the LLM extraction.
First, the provenance script adds \texttt{source\_site}, \texttt{source\_url},
\texttt{original\_title}, and \texttt{collection\_date} to every available JSON
record using collection indexes and archived source files. These values are not
generated by the LLM. The script reports coverage by source and leaves a field
empty when a reliable value is unavailable. In particular, the SIAGE action
sheets are marked as internal material and do not receive an artificial public
URL. Historical collection dates are also left empty because they were not
consistently recorded in the original indexes.

Second, a taxonomy-normalisation script adds
\texttt{thematique\_normalisee} and
\texttt{nature\_initiative\_normalisee}, while preserving the original LLM
fields. The implementation maps lexical variants to a documented French
controlled taxonomy so that related labels can be aggregated consistently in
filters and charts. It does not invoke the LLM again. Ambiguous values are
assigned to \textit{autre / à revoir} and written to a review report instead of
being forced into an unsuitable category. The records also store a taxonomy
version, allowing later revisions to be identified and reapplied.

\section{Main modules and outputs}

The main modules can be grouped into five functional families:

\begin{itemize}
    \item source-specific collection scripts for CNAV, Fondation de France,
    HCFEA, IReSP, RFVAA, and the SIAGE action sheets;
    \item cleaning, file-name standardisation, segmentation, and
    keyword-based relevance-scoring scripts;
    \item rule-based extraction modules that generate intermediate candidates,
    scores, consolidated records, and review tables;
    \item local LLM extraction, JSON repair, taxonomy normalisation, and
    provenance-enrichment scripts;
    \item evaluation and analysis scripts that create reference samples,
    calculate scores, and produce spreadsheets and charts.
\end{itemize}

The principal output is a collection of JSON records usable by downstream
software. Complementary outputs include a consolidated JSON file, CSV review
tables, Excel workbooks, and charts describing themes, initiative types,
stakeholders, territories, temporal mentions, and comparisons between SIAGE
action sheets and other documents. The current analytical scripts prioritise
the normalised theme and initiative-type fields when available, while retaining
the original values for audit and correction. These outputs are exploratory;
they do not estimate the real-world frequency or effectiveness of initiatives
in France.

\section{Difficulties and implemented solutions}

Structural heterogeneity was the main implementation difficulty: a heading in
one source can be expressed as a paragraph, table, caption, or be absent in
another. The modular pipeline and the separation between raw, cleaned, and
intermediate directories made this variability inspectable. The transition from
rule-based extraction to constrained LLM extraction reduced the need for a
large number of source-specific rules, while controlled vocabularies and
post-extraction normalisation limited avoidable output variation.

Several difficulties remain. PDF conversion can introduce encoding artefacts or
reading-order errors, and websites may change their structure after collection.
Dates can refer to publication, funding, event, or initiative-start dates;
territorial references can refer to a place, an organisation, or a project
area; and an organisation mentioned in a document is not necessarily the
initiative holder. The system therefore retains source links, original titles,
raw texts, and intermediate outputs so that these cases can be reviewed. The
human reference evaluation described in the next chapter provides a separate
assessment of factual extraction quality; the implementation itself does not
treat field completion as evidence of correctness.
